\section{Introduction}\label{sec:introduction}

The spatial queries that dominate point-cloud processing are not the sparse requests of an interactive user but the dense, machine-generated ones of a pipeline: normal estimation asks for the nearest neighbours of every point, outlier removal and clustering ask for the points within a radius of every point, and feature computation does both again at several scales. A cloud of a hundred million points puts hundreds of millions of neighbourhood queries through its index per stage, and surveys of a billion points are now routine. The index is the inner loop of the application, and its choice is a design decision that a practitioner makes once, by habit or by a library's default.

The decision has several parts. The first is the primitive: a kd-tree, an octree, a quadtree, a grid, an R-tree or a bounding-volume hierarchy, each with its own fan-out and its own way of cutting space. The second is the parameters: the leaf capacity, the depth, the split rule. The third is the memory layout: how nodes and points are stored and visited. A fourth part has been proposed repeatedly over the last decade: \emph{composition}, the nesting of several primitives in one hierarchy so that the structure changes with depth, a wide one at the top and a fine one below. Several hybrid indexes for point clouds have been reported to beat single structures~\cite{wang2021hybridpointcloud,park2021nestedoctree,ogayar2023nestedlidar}, among them the nested data structures for LiDAR of Ogayar-Anguita et al.~\cite{ogayar2023nestedlidar}, whose quadtree-over-octree schemas loaded fewer nodes than either structure alone on airborne scans. The evidence for composition has two recurring features: it is measured against single structures at default parameters, and it is counted in nodes visited or loaded rather than in time.

This paper asks a simpler question than the hybrid literature does. When every candidate, single or nested, is built by the same engine into the same memory layout, tuned under the same budget, and measured on the queries a real pipeline issues, held out from the selection, where does the speed come from? And does composition survive that measurement?

To answer it we built an instrument rather than an index: an engine that compiles any composition of primitives into one node pool and one traversal, a grammar of schemas that makes a single structure the one-block case of a nested one, and a search that is free to return either. We ran it on \nCells{} cells of \nScenes{} scenes from one million to one billion points, from terrestrial and mobile scans to airborne surveys and forests, with query traces recorded from three pipeline stages and split into a half the search sees and a half it is judged on. We then took the question out of memory, to seven other workloads, to the GPU, to emulated node costs, and to out-of-core stores with the leaves and then the hierarchy on four classes of device; we timed five external libraries on the same traces under one protocol; and we priced every node visit by depth to see where a query's time goes.

The findings are these.
\begin{itemize}[leftmargin=*, itemsep=2pt]
\item \textbf{The primitive and the leaf are the levers.} Choosing the primitive and tuning its leaf, depth and split rule account for the whole gain between a library default and the best structure a search can find, a factor of \leverOctreeDefaultTuned{} over an octree default and \leverKdDefaultTuned{} over a median kd-tree, at every size from one million to one billion points. Composition adds nothing: the best nested schema is \nestedSingleAll$\times$ the best single over \nCells{} cells, was never selected, and the picture is the same on seven workloads, on the GPU, under emulated node costs, and with the leaves out of core.
\item \textbf{Composition pays only with the hierarchy paged.} A nested prefix is adopted in \oocPagedAdopted{} of \oocPagedCells{} out-of-core cells, all with the hierarchy itself read page by page from a slow device, which is the regime of the 2023 result and accounts for it; pinning the hierarchy, a few tens of megabytes at 100M points, is several times faster than any paged schema.
\item \textbf{One feature picks the primitive.} A single split on a scene's flatness score, kd-tree below and quadtree above, is within \ruleRegret{} to \ruleRegretTimed$\times$ of a per-cloud search, closing \ruleCapturePct\% of the gap from the default; nanoflann's default is already that structure for three-dimensional scenes and ties the tuned pick, while four other libraries are 1.6 to 2.6$\times$ slower by implementation.
\item \textbf{Node counts do not predict time; layout does.} A node visit costs \visitCostTopNs\,ns near the root and \visitCostLeafNs\,ns at a leaf, so the visits a wide prefix saves are the cheap ones, and count-based cost models explain under half of the measured lever. Two-megabyte pages speed every structure by \hugeSinglesHundredM$\times$ at 100M points and up to \hugeMin$\times$ for the largest pools, more than composition ever did, and they move the optimal leaf, so a result measured on 4\,KiB pages carries the page size with it.
\end{itemize}

The paper is a revisit, and we say so: one of its findings reverses a result published by some of the present authors' community, and the mechanism that explains the reversal, untuned baselines and node counts in place of time, is the mechanism of the hybrid-index literature at large. What we offer in its place is smaller and, we think, more useful: a controlled protocol, a released instrument with its traces, configurations and scripts, and a rule a practitioner can apply without a search.

Section~\ref{sec:related} places the work. Section~\ref{sec:instrument} describes the engine, the grammar and the search; Section~\ref{sec:setup} the scenes, workloads, protocol and hardware. Section~\ref{sec:results} reports the levers, composition in and out of memory, the rule, the libraries and the mechanism. Section~\ref{sec:discussion} discusses why hybrids have looked good, what transfers, and what the study does not settle.
